% Options for packages loaded elsewhere
\PassOptionsToPackage{unicode}{hyperref}
\PassOptionsToPackage{hyphens}{url}
\documentclass[
  a4paper,
]{ltjsarticle}
\usepackage{xcolor}
\usepackage[margin=25mm]{geometry}
\usepackage{amsmath,amssymb}
\setcounter{secnumdepth}{-\maxdimen} % remove section numbering
\usepackage{iftex}
\ifPDFTeX
  \usepackage[T1]{fontenc}
  \usepackage[utf8]{inputenc}
  \usepackage{textcomp} % provide euro and other symbols
\else % if luatex or xetex
  \usepackage{unicode-math} % this also loads fontspec
  \defaultfontfeatures{Scale=MatchLowercase}
  \defaultfontfeatures[\rmfamily]{Ligatures=TeX,Scale=1}
\fi
\usepackage{lmodern}
\ifPDFTeX\else
  % xetex/luatex font selection
\fi
% Use upquote if available, for straight quotes in verbatim environments
\IfFileExists{upquote.sty}{\usepackage{upquote}}{}
\IfFileExists{microtype.sty}{% use microtype if available
  \usepackage[]{microtype}
  \UseMicrotypeSet[protrusion]{basicmath} % disable protrusion for tt fonts
}{}
\makeatletter
\@ifundefined{KOMAClassName}{% if non-KOMA class
  \IfFileExists{parskip.sty}{%
    \usepackage{parskip}
  }{% else
    \setlength{\parindent}{0pt}
    \setlength{\parskip}{6pt plus 2pt minus 1pt}}
}{% if KOMA class
  \KOMAoptions{parskip=half}}
\makeatother
\usepackage{longtable,booktabs,array}
\newcounter{none} % for unnumbered tables
\usepackage{calc} % for calculating minipage widths
% Correct order of tables after \paragraph or \subparagraph
\usepackage{etoolbox}
\makeatletter
\patchcmd\longtable{\par}{\if@noskipsec\mbox{}\fi\par}{}{}
\makeatother
% Allow footnotes in longtable head/foot
\IfFileExists{footnotehyper.sty}{\usepackage{footnotehyper}}{\usepackage{footnote}}
\makesavenoteenv{longtable}
\usepackage{graphicx}
\makeatletter
\newsavebox\pandoc@box
\newcommand*\pandocbounded[1]{% scales image to fit in text height/width
  \sbox\pandoc@box{#1}%
  \Gscale@div\@tempa{\textheight}{\dimexpr\ht\pandoc@box+\dp\pandoc@box\relax}%
  \Gscale@div\@tempb{\linewidth}{\wd\pandoc@box}%
  \ifdim\@tempb\p@<\@tempa\p@\let\@tempa\@tempb\fi% select the smaller of both
  \ifdim\@tempa\p@<\p@\scalebox{\@tempa}{\usebox\pandoc@box}%
  \else\usebox{\pandoc@box}%
  \fi%
}
% Set default figure placement to htbp
\def\fps@figure{htbp}
\makeatother
\setlength{\emergencystretch}{3em} % prevent overfull lines
\providecommand{\tightlist}{%
  \setlength{\itemsep}{0pt}\setlength{\parskip}{0pt}}
% Glyph map for lualatex (newunicodechar). Maps symbols that may lack font glyphs.
\usepackage{newunicodechar}
\usepackage{amssymb}
\newunicodechar{∎}{\ensuremath{\blacksquare}}
\newunicodechar{✓}{\ensuremath{\checkmark}}
\newunicodechar{✗}{\ensuremath{\times}}
\newunicodechar{⟹}{\ensuremath{\Longrightarrow}}
\newunicodechar{⟺}{\ensuremath{\Longleftrightarrow}}
\newunicodechar{↪}{\ensuremath{\hookrightarrow}}
\newunicodechar{⌊}{\ensuremath{\lfloor}}
\newunicodechar{⌋}{\ensuremath{\rfloor}}
\newunicodechar{≃}{\ensuremath{\simeq}}
\newunicodechar{≈}{\ensuremath{\approx}}
\newunicodechar{≤}{\ensuremath{\leq}}
\newunicodechar{≥}{\ensuremath{\geq}}
\newunicodechar{≠}{\ensuremath{\neq}}
\newunicodechar{→}{\ensuremath{\rightarrow}}
\newunicodechar{←}{\ensuremath{\leftarrow}}
\newunicodechar{↔}{\ensuremath{\leftrightarrow}}
\newunicodechar{⊥}{\ensuremath{\perp}}
\newunicodechar{√}{\ensuremath{\surd}}
\newunicodechar{∑}{\ensuremath{\sum}}
\newunicodechar{∏}{\ensuremath{\prod}}
\newunicodechar{∞}{\ensuremath{\infty}}
\newunicodechar{∈}{\ensuremath{\in}}
\newunicodechar{∀}{\ensuremath{\forall}}
\newunicodechar{∣}{\ensuremath{\mid}}
\newunicodechar{γ}{\ensuremath{\gamma}}
\newunicodechar{λ}{\ensuremath{\lambda}}
\newunicodechar{μ}{\ensuremath{\mu}}
\newunicodechar{ρ}{\ensuremath{\rho}}
\newunicodechar{σ}{\ensuremath{\sigma}}
\newunicodechar{φ}{\ensuremath{\varphi}}
\newunicodechar{ψ}{\ensuremath{\psi}}
\newunicodechar{θ}{\ensuremath{\theta}}
\newunicodechar{Δ}{\ensuremath{\Delta}}
\newunicodechar{Π}{\ensuremath{\Pi}}
\newunicodechar{Λ}{\ensuremath{\Lambda}}
\newunicodechar{τ}{\ensuremath{\tau}}
\newunicodechar{ε}{\ensuremath{\varepsilon}}
\newunicodechar{ω}{\ensuremath{\omega}}
\newunicodechar{π}{\ensuremath{\pi}}
\newunicodechar{ν}{\ensuremath{\nu}}
\newunicodechar{±}{\ensuremath{\pm}}
\newunicodechar{×}{\ensuremath{\times}}
\newunicodechar{·}{\ensuremath{\cdot}}
\usepackage{bookmark}
\IfFileExists{xurl.sty}{\usepackage{xurl}}{} % add URL line breaks if available
\urlstyle{same}
\hypersetup{
  pdftitle={関係波閉鎖系における単一波の運動方程式------L(K\_N) 上の蔵本型位相結合と、90°位相差骨格での自己無撞着床の条件},
  pdfauthor={Noriaki Kihara (WF System Co., Ltd.), ORCID 0009-0004-6753-4020},
  hidelinks,
  pdfcreator={LaTeX via pandoc}}

\title{関係波閉鎖系における単一波の運動方程式------\(L(K_N)\)
上の蔵本型位相結合と、90°位相差骨格での自己無撞着床の条件}
\author{Noriaki Kihara (WF System Co., Ltd.), ORCID 0009-0004-6753-4020}
\date{2026-09-12}

\begin{document}
\maketitle

\textbf{シリーズ:} 自己無撞着な関係波閉鎖系の基礎（論文1／全10本）\\
\textbf{著者:} 木原 範昭（WF System Co., Ltd.）　\textbf{日付:}
2026-09-12\\
\textbf{Version DOI:} 10.5281/zenodo.22728761\\
\textbf{Concept DOI:} 10.5281/zenodo.22728760\\
\textbf{ORCID:} 0009-0004-6753-4020\\
\textbf{Zenodo:} https://zenodo.org/records/22728761\\

\begin{center}\rule{0.5\linewidth}{0.5pt}\end{center}

\subsection{要旨}\label{ux8981ux65e8}

完全グラフ \(K_N\) の各辺に複素関係波 \(z_e\in\mathbb C\)
を置き位相のみの実反対称生成子で回す閉鎖系写像（先行研究［1］で定義）について、本論文の問いは経験的動機から立つ：乱数シードから自己無撞着固定点
\(KZ=i\sigma Z\) を作る make\_parent の床が全 \(N\)
で\textbf{自然に90°二軸位相格子}になった（実験4-1・検討28）ため、『なぜ90°か』ではなく------90°を所与として------『\textbf{90°位相差骨格の上に、自己無撞着に回転する極めて対称な床を構成できるか}』を問う。まず
one\_step 写像 \(z\to\exp(\Delta\tau K)z\)
を振幅・位相に分解し、各辺波が頂点共有の隣接辺とのみ蔵本型
\(\sin(\varphi_f-\varphi_e)\) で結合することと、

\[\frac{dr_e}{d\tau}=\frac12\sum_{f\sim e}r_f\sin\!\big(2(\varphi_f-\varphi_e)\big)\quad(\text{振幅}),\qquad r_e\frac{d\varphi_e}{d\tau}=\sum_{f\sim e}r_f\sin^2(\varphi_f-\varphi_e)\quad(\text{位相})\]

を導く。90°骨格（位相差が \(90^\circ\) の倍数）では振幅駆動の第2高調波
\(\sin(2\psi)\) が \textbf{termwise に消え}
\(\dot r_e=0\)（振幅凍結）。残る自己無撞着条件（一様回転
\(\dot\varphi_e=\omega\)）は振幅ベクトルの固有値問題
\(\omega r_e=\sum_{f\sim e}r_f\sin^2\psi_{ef}\)（\(Wr=\omega r\)）に帰着する。さらに
\(D=\operatorname{diag}(e^{i\varphi_e})\) に対し90°骨格では
\(\boxed{D^{-1}KD=iW}\)（\(W_{ef}=A_{ef}\sin^2\psi_{ef}\)
実対称・非負）が成り立ち、\textbf{複素自己無撞着問題
\(Kz=i\omega z\)（\(z=Dr\)）は実対称非負整数行列の固有値問題
\(Wr=\omega r\) と厳密同値}である。\(W\) 既約なら Perron--Frobenius
で正振幅の主固有ベクトルが一意に存在し（\(\omega=\rho(W)\)）、\(z=Dr\)
が90°骨格上の自己無撞着波を構成する（全 \(N=3\)--\(40\) で \(W\)
既約・確認）。その代数的厳密解は論文3が与える（make\_parent
床は各自の90°ラベルが定める \(W\) の Perron
固有ベクトルの反復近似で、整数K床とはラベル・\(W\)
が異なる別メンバー）。同一Nで並べると、make\_parent 床の振幅は \(N=40\)
で699値とほぼ全辺で異なるのに対し、意図構成の高対称床は振幅≤3値で対称性が明確に高い（図2）。厳密写像の軌道再現は
\(\le3.0\times10^{-15}\)、EOMの有限差分検証は振幅
\(\le1.4\times10^{-6}\)・位相
\(\le2.4\times10^{-7}\)、90°骨格の振幅凍結は
\(\le1.1\times10^{-13}\)（\(N=6,22,40\)）。分類：正本 one\_step
からの導出済み帰結＋機械精度検証。

\begin{center}\rule{0.5\linewidth}{0.5pt}\end{center}

\subsection{1. はじめに}\label{ux306fux3058ux3081ux306b}

本シリーズは、\(N\) 個の存在の全二体関係 \(M=N(N-1)/2\) を複素関係波
\(z_e\) とし、外部seedを除去した閉鎖系の力学を調べる。初期に用いた
\textbf{make\_parent}
は、乱数シードから自己無撞着な固定点（\(KZ=i\sigma Z\)）を反復で作る手続きで、その床は全
\(N=3\)--\(40\) で\textbf{自然に90°二軸位相格子}（相対位相差
\(90^\circ\)）になった（図1、実験4-1・検討28）。位相は自然に90°だが、振幅は辺ごとにバラバラ（\(N=40\)
で699値・CV数\%）・重心 \(\Sigma z\neq0\) で、対称性は高くない。

\begin{figure}
\centering
\pandocbounded{\includegraphics[keepaspectratio,alt={図1: 乱数make\_parent床（乱数シードからの自己無撞着固定点 KZ=iσZ）の複素平面（全N=3-40, step0）。全Nで全波が直交2軸（相対位相差90°）に乗る＝自然に90°二軸格子になる。振幅は辺ごとにバラバラ。}]{figs/p1_single_wave_eom_ja_fig1.png}}
\caption{図1: 乱数make\_parent床（乱数シードからの自己無撞着固定点
KZ=iσZ）の複素平面（全N=3-40,
step0）。全Nで全波が直交2軸（相対位相差90°）に乗る＝自然に90°二軸格子になる。振幅は辺ごとにバラバラ。}
\end{figure}

\textbf{図1} 乱数 make\_parent
床の複素平面（\(N=3\)--\(40\)、step0）。乱数シードから作った自己無撞着固定点が、全
\(N\) で自然に直交2軸（相対位相差
\(90^\circ\)）に乗る。これが本論文の出発点＝「90°は所与」の経験的事実であり、その理由（自己無撞着＋termwise
から90°が必然）は論文0が与える。

この経験的事実を出発点に、本論文の問いは『なぜ90°か』ではなく------\textbf{90°は所与（自然にそうなった）とした上で}------『\textbf{90°位相差骨格の中に、自己無撞着に回転する極めて対称な床を構成できるか}』である。実際、同一
\(N\) で並べると（図2）、乱数 make\_parent 床（振幅ほぼ \(M\)
値）に対し、意図的に構成した高対称床は振幅が少数値（高対称v2：≤2値・\(D_N\)対称／整数K解析床：≤3値・小重心の相対平衡）で、対称性が明確に高い。

本論文はまず単一波EOMを導出し、90°骨格で振幅が termwise
に凍ること、残る自己無撞着条件が振幅ベクトルの固有値問題になること（→論文3で厳密解）を示す。床の点火資格・厳密解・対称構成はそれぞれ論文2・3で扱う。力学の定義以外の入力（零閉塞・特定振幅分布）は前提しない。

\begin{figure}
\centering
\pandocbounded{\includegraphics[keepaspectratio,alt={図2: 床の対称性比較（同一N、位相はどれも90°）。左=乱数make\_parent床（振幅ほぼM値・重心≠0）、中=高対称v2床（振幅≤2値だが重心最大級）、右=整数K解析床（≤3値・小重心・相対平衡＝論文3）。}]{figs/p1_single_wave_eom_ja_fig2.png}}
\caption{図2:
床の対称性比較（同一N、位相はどれも90°）。左=乱数make\_parent床（振幅ほぼM値・重心≠0）、中=高対称v2床（振幅≤2値だが重心最大級）、右=整数K解析床（≤3値・小重心・相対平衡＝論文3）。}
\end{figure}

\textbf{図2} 床の対称性比較（\(N=6,12,17,40\)）。乱数 make\_parent
が自然に生む90°床（振幅ほぼ \(M\) 値・\(N=40\)
で699値）に対し、意図構成の高対称床は振幅が少数値で対称性が高い（詳細は論文2・3）。

\subsection{2. 力学（写像）}\label{ux529bux5b66ux5199ux50cf}

正本 \texttt{run\_N3\_N40\_stage123\_v1.py}（SHA256
\texttt{1abf2353…}）の one\_step
は、\(u=e^{i\arg z}\)、\(H_{ef}=A_{ef}\,u_f\bar u_e\)、\(H\leftarrow i\,\Im H\)
として \(z_{n+1}=\exp(-i\Delta\tau H)z_n\)。整理すると

\[K_{ef}=A_{ef}\sin(\varphi_f-\varphi_e),\qquad z_{n+1}=\exp(\Delta\tau K)\,z_n,\qquad \Delta\tau=\frac{2\pi}{N},\]

ここで \(A\) は \(K_N\) の辺グラフ \(L(K_N)\) の隣接行列（辺 \(e,f\)
が頂点共有なら \(A_{ef}=1\)）、\(\varphi_e=\arg z_e\)。\textbf{\(K\)
は位相のみで決まり振幅を含まない実反対称行列}で、\(\exp(\Delta\tau K)\)
は実直交行列。\(\Re z,\Im z\) に同一作用するため \(\|z\|^2\) と二乗閉塞
\(z^{\mathsf T}z=\sum_e z_e^2\) を厳密保存する。

\subsection{3.
単一波の運動方程式（導出）}\label{ux5358ux4e00ux6ce2ux306eux904bux52d5ux65b9ux7a0bux5f0fux5c0eux51fa}

生成子の流れ \(dz/d\tau=Kz\)
を成分で書くと、\(f\sim e\)（頂点共有辺）について

\[\frac{dz_e}{d\tau}=\sum_{f\sim e}\sin(\varphi_f-\varphi_e)\,z_f=\sum_{f\sim e}r_f\sin(\varphi_f-\varphi_e)\,e^{i\varphi_f}.\]

\(z_e=r_e e^{i\varphi_e}\) とし左辺を
\(e^{i\varphi_e}(\dot r_e+ir_e\dot\varphi_e)\)、右辺の
\(e^{i\varphi_f}=e^{i\varphi_e}e^{i\psi}\)（\(\psi=\varphi_f-\varphi_e\)）で
\(e^{i\varphi_e}\) を約すと
\(\text{RHS}=\sum r_f\sin\psi(\cos\psi+i\sin\psi)\)。実部・虚部より

\[\boxed{\ \dot r_e=\frac12\sum_{f\sim e}r_f\sin(2\psi)\ }\qquad\boxed{\ r_e\dot\varphi_e=\sum_{f\sim e}r_f\sin^2\psi\ }\]

を得る（\(\sin\psi\cos\psi=\tfrac12\sin2\psi\)
を用いた）。各辺波は頂点共有の隣接辺とのみ蔵本型に結合する振動子である。

\subsection{4.
この式が軌道を説明する}\label{ux3053ux306eux5f0fux304cux8eccux9053ux3092ux8aacux660eux3059ux308b}

\begin{itemize}
\tightlist
\item
  \textbf{位相（回転と剛体回転を除いた変位）}：\(\sin^2\ge0\) ゆえ全波
  \(\varphi_e\) は単調前進＝回転。共通分が剛体回転、波ごとの
  \((1/r_e)\sum r_f\sin^2\psi\)
  の\textbf{差}が剛体回転を除いた位相変位。
\item
  \textbf{振幅（増減）}：駆動は第2高調波
  \(\sin(2\psi)\)。隣接辺との位相差の2倍の正弦の重み付き和。
\item
  \textbf{床（90°骨格）}：位相差が全て \(90^\circ\) の倍数
  \(\Rightarrow\sin(2\psi)=0\Rightarrow\dot r_e=0\)（termwise
  振幅凍結）。さらに振幅が
  \(Wr=\omega r\)（\(W_{ef}=A_{ef}\sin^2\psi_{ef}\)）を満たす配置では全波が同率回転
  \(\dot\varphi_e=\omega\)
  となり、自己無撞着相対平衡を形成する。この床の存在・一意性は §5
  で扱う。
\end{itemize}

本論文の主題は90°骨格上の自己無撞着床の\textbf{構成}であり、その床が不安定であること（点火の資格）と、床を離れた後の動態（インフレーション・再ロック）は本論文では扱わない（論文2・5・6）。

\subsection{5.
90°骨格での振幅凍結と自己無撞着条件}\label{ux9aa8ux683cux3067ux306eux632fux5e45ux51cdux7d50ux3068ux81eaux5df1ux7121ux649eux7740ux6761ux4ef6}

90°は所与の骨格（make\_parent
の自己無撞着床が自然にそうなった、§1）である。その上で本節は、90°骨格が振幅を
termwise
に凍らせること、そして自己無撞着床の条件が振幅の固有値問題になることを示す。

\textbf{振幅凍結（termwise）}：振幅駆動は第2高調波 \(\sin(2\psi)\)
で、\([0,360^\circ)\) の零点は \(\psi=0,90,180,270^\circ\)
のみ（数え上げ）。ゆえに位相差が全て \(90^\circ\) の倍数なら全対で
\(\sin(2\psi)=0\)、すなわち \(\dot r_e=0\)
が\textbf{各項ごと}に成り立つ（termwise）。これは生成子が
\(\sin\psi\)（第1高調波）ゆえ振幅駆動が第2高調波になる直接の帰結で、90°骨格が「振幅を壊さず置ける」骨格である理由である。（一般位相でも辺ごとの和が相殺すれば
\(\dot r=0\) になり得る＝cancellation
型で、これはロック終点に対応する。）

\textbf{定理（90°骨格の実対称化）}
90°骨格の上で自己無撞着な単一回転（\(\dot\varphi_e=\omega\)
一様）を課すと、位相方程式は
\(\omega r_e=\sum_{f\sim e}r_f\sin^2\psi_{ef}\)、すなわち
\(Wr=\omega r\)（\(W_{ef}=A_{ef}\sin^2\psi_{ef}\)）になる。これは固有値問題が単に「出てくる」以上の\textbf{厳密同値}である：\(\psi_{ef}\in\frac{\pi}{2}\mathbb Z\)、\(D=\operatorname{diag}(e^{i\varphi_e})\)
なら、恒等式
\(\sin\psi\,e^{i\psi}=\tfrac12\sin2\psi+i\sin^2\psi=i\sin^2\psi\)（\(\Leftrightarrow\sin2\psi=0\)）より
\[\boxed{\ D^{-1}KD=iW\ },\qquad W_{ef}=A_{ef}\sin^2\psi_{ef}\in\{0,1\}\ (\text{実対称・非負}).\]
したがって
\(\boxed{\,Kz=i\omega z\ (z=Dr,\ r\ \text{実})\iff Wr=\omega r\,}\)。90°骨格の複素自己無撞着固有波問題は、実対称非負整数行列
\(W\) の固有値問題と厳密に同値である。

\textbf{系（正振幅自己無撞着波の存在・一意性）} \(W\)
が既約（対応グラフが連結）なら、Perron--Frobenius により
\(Wr=\rho(W)r\)、\(r_e>0\) を満たす \(r\)
がスケールを除き一意に存在する。よって \(\boxed{\,z=Dr\,}\)
は90°骨格上の自己無撞着単一回転 \(Kz=i\rho(W)z\)
である。すなわち本論文の問い「90°骨格上に自己無撞着な正振幅波を作れるか」の答えは、\(W\)
既約の下で \textbf{作れる（スケールを除き一意）}。

\textbf{確認（全 \(N=3\)--\(40\)・追試）} 整数K床・make\_parent 床とも
\(\|D^{-1}KD-iW\|\le2\times10^{-12}\)、\(Wr=\omega r\) の
\(\omega=\rho(W)=\sigma\)（iK スペクトル半径）一致、\(r>0\)、\(r\) と
\(W\) 主固有ベクトルの重なり \(1.000000\)、\(W\) は全 \(N\)
で既約（連結）。要約すれば
\(\boxed{\text{90°骨格}\overset{D^{-1}KD=iW}{\longrightarrow}\text{実非負対称固有値問題}\overset{\rm PF}{\longrightarrow}\text{正振幅自己無撞着波の存在・一意}}\)。

その高対称解の\textbf{閉形式}（\(\rho(W)\) の \(N\)
依存、振幅の2/3値縮退、整数 \(K\) の \(\sigma_{\max}\)
との一致）は論文3が具体的に解く。make\_parent
床は、反復の結果として得られた各自の90°位相ラベルが定める
\(W\)（\(W_{\rm mp}\)）の Perron
主固有ベクトルを近似したものであり（和則
\((\sum_f r_f\sin^2\psi_{ef})/r_e=\sigma\) が全 \(N\)
で成立、検討31）、整数K床とはラベルも \(W\)
も異なる別メンバーだが、各々が自分の \(W\) の Perron
解である点で整合する（\(W_{\rm mp}\neq W_{\rm integerK}\)、論文0）。

すなわち90°は導出される必然でなく所与の骨格であり、本論文が確定するのは「90°骨格では振幅が
termwise に凍り、自己無撞着床の存在が振幅の固有値問題 \(Wr=\omega r\)
に帰着する」ことである。床の点火資格（不安定性）は論文2、厳密解は論文3、そこからの急拡大は論文5・6で扱う。

\subsection{\texorpdfstring{6.
数値検証（\(N=6,22,40\)）}{6. 数値検証（N=6,22,40）}}\label{ux6570ux5024ux691cux8a3cn62240}

{\def\LTcaptype{none} % do not increment counter
\begin{longtable}[]{@{}
  >{\raggedright\arraybackslash}p{(\linewidth - 4\tabcolsep) * \real{0.3333}}
  >{\raggedright\arraybackslash}p{(\linewidth - 4\tabcolsep) * \real{0.3333}}
  >{\raggedright\arraybackslash}p{(\linewidth - 4\tabcolsep) * \real{0.3333}}@{}}
\toprule\noalign{}
\begin{minipage}[b]{\linewidth}\raggedright
テスト
\end{minipage} & \begin{minipage}[b]{\linewidth}\raggedright
内容
\end{minipage} & \begin{minipage}[b]{\linewidth}\raggedright
最大誤差
\end{minipage} \\
\midrule\noalign{}
\endhead
\bottomrule\noalign{}
\endlastfoot
T1 & 厳密写像 \(\exp(\Delta\tau K(\varphi_n))z_n\) が正本軌道
\(z_{n+1}\) を再現（\(K\) が生成子） & \(\le3.0\times10^{-15}\) \\
T2 & サブステップ \((\exp(\varepsilon K)z-z)/\varepsilon\)
の振幅・位相成分が §3 のEOMと一致 & 振幅 \(\le1.4\times10^{-6}\)／位相
\(\le2.4\times10^{-7}\) \\
T3 & 床の90°格子で \(\dot r_e\approx0\)（振幅凍結） &
\(|\dot r|\le1.1\times10^{-13}\) \\
\end{longtable}
}

（T2は \(\varepsilon\) スケール差分の丸め由来。式は厳密。）90°骨格 vs
ロック終点の実測：90°骨格（step0）は個別項
\(|\sin2\psi|\sim10^{-12}\text{–}10^{-14}\)（termwise、格子ずれ
\(0.0^\circ\)）、ロック終点は個別項 \(\sim1.0\) だが辺ごと重み和
\(\sim10^{-7}\text{–}10^{-13}\)（cancellation、格子ずれ
\(\sim45^\circ\)）。

\begin{figure}
\centering
\pandocbounded{\includegraphics[keepaspectratio,alt={図3: 単一波EOMの検証（N=6）。左=振幅、右=位相、実測（縦）と導出式（横）が y=x 上に一致。}]{figs/p1_single_wave_eom_ja_fig3.png}}
\caption{図3:
単一波EOMの検証（N=6）。左=振幅、右=位相、実測（縦）と導出式（横）が y=x
上に一致。}
\end{figure}

\textbf{図3} 単一波運動方程式の検証（\(N=6\)）。

\begin{figure}
\centering
\pandocbounded{\includegraphics[keepaspectratio,alt={図4: 90°骨格の振幅凍結（termwise）とロック（cancellation）。左=sin(2ψ)の零点=k·90°、右=90°骨格（termwise, 個別項\textasciitilde0）とロック（cancellation, 個別項O(1)）の実測分布（N=6）。}]{figs/p1_single_wave_eom_ja_fig4.png}}
\caption{図4:
90°骨格の振幅凍結（termwise）とロック（cancellation）。左=sin(2ψ)の零点=k·90°、右=90°骨格（termwise,
個別項\textasciitilde0）とロック（cancellation,
個別項O(1)）の実測分布（N=6）。}
\end{figure}

\textbf{図4} 90°骨格（termwise
振幅凍結）とロック終点（cancellation）の実測（\(N=6\)）。

\subsection{7.
主張分類・未決}\label{ux4e3bux5f35ux5206ux985eux672aux6c7a}

\begin{itemize}
\tightlist
\item
  分類：正本 one\_step
  からの\textbf{導出済み帰結}（＋機械精度検証）。判定：保持。
\item
  断定（本論文の主結果）：90°骨格で \(D^{-1}KD=iW\)
  が成り立ち、複素自己無撞着問題 \(Kz=i\omega z\) は実対称非負整数行列
  \(W_{ef}=A_{ef}\sin^2\psi_{ef}\) の固有値問題 \(Wr=\omega r\)
  と\textbf{厳密同値}。\(W\) 既約なら Perron--Frobenius
  で正振幅の主固有ベクトルが一意（\(\omega=\rho(W)\)）＝\textbf{90°骨格上に自己無撞着な正振幅床が構成できる}（全
  \(N=3\)--\(40\) で \(W\)
  既約・機械精度確認）。これが「90°骨格に自己無撞着な高対称波を作れるか」への数学的回答である。
\item
  未決：床の全体スケール \(r^2\)（スケール私有）、相対平衡の
  moduli（90°床は最も対称な一員で唯一解でない）、\(N=3\)
  の例外（\(\mu/r^2=-3/2\), rank 1）、任意の90°ラベル配置での \(W\)
  既約性の一般条件。これらは論文2・3・5で扱う。
\end{itemize}

\subsection{関連研究（構造的一致・導出元でない）}\label{ux95a2ux9023ux7814ux7a76ux69cbux9020ux7684ux4e00ux81f4ux5c0eux51faux5143ux3067ux306aux3044}

\begin{itemize}
\tightlist
\item
  蔵本モデル（Kuramoto
  1975/1984）：\(K_{ef}=A_{ef}\sin(\varphi_f-\varphi_e)\) は \(L(K_N)\)
  上の蔵本型位相結合。
\end{itemize}

\subsection{参考文献}\label{ux53c2ux8003ux6587ux732e}

\begin{enumerate}
\def\labelenumi{\arabic{enumi}.}
\tightlist
\item
  木原範昭「自己無撞着な関係波閉鎖系におけるインフレーション的急拡大の機構」Concept
  DOI 10.5281/zenodo.22112008.（本論文の力学の定義・\(\tau\neq\)
  時間の出典）
\item
  Y. Kuramoto, \emph{Chemical Oscillations, Waves, and Turbulence},
  Springer (1984); ``Self-entrainment of a population of coupled
  non-linear oscillators'', 1975.
\end{enumerate}

\subsection{再現}\label{ux518dux73fe}

検証スクリプト・データは
\texttt{位相ロック全N確認\_相対平衡\_20260912/単一波運動方程式\_導出と検証\_20260912/}（\texttt{verify\_single\_wave\_eom\_20260912.py}＝T1/T2/T3、\texttt{verify\_90deg\_necessity\_20260912.py}、\texttt{make\_figures\_eom\_90deg\_20260912.py}、\texttt{REPORT.md}、\texttt{SHA256SUMS.txt}、\texttt{run\_all.sh}）に一式。\textbf{90°骨格での複素↔実対称同値
\(D^{-1}KD=iW\) と Perron--Frobenius の全 \(N\) 検証（§5）}は
\texttt{位相ロック全N確認\_相対平衡\_20260912/複素実対称同値\_DKD\_iW\_20260912/}（\texttt{verify\_DKD\_iW\_equivalence\_20260912.py}、\texttt{results/DKD\_iW\_summary.csv}、README／REPORT／SHA／run\_all；読出しのみ）。図1の対称性比較は
\texttt{同/床の対称性比較\_makeparent\_vs\_高対称\_20260912/}。入力は既存
den=N npz（\(N=6\): 500 step、\(N=22,40\): 本系列 2000
step）。力学写像は正本 \texttt{run\_N3\_N40\_stage123\_v1.py}（SHA256
\texttt{1abf2353fee2e4f56f05e7a6f149fd086885136beb61ab571b48a56b09691567}）と同一・物理式無変更。

\end{document}
